\chapter{Implementation}\label{ch:implementation}

\section{Software, Hardware and Application}

The system is written in Python with PyTorch \cite{r22}, TransformerLens 3.5.1 \cite{r13} (named hook points on every GPT-2 activation, and running the model from a saved intermediate state), SAELens 6.45.3 \cite{r14} (loading the \texttt{gpt2-small-res-jb} dictionaries \cite{r15}), Neuronpedia \cite{r20} for feature descriptions and Streamlit for the application. The main-line studies ran on a GTX 1660 Ti PC (PyTorch 2.6.0, CUDA 12.4); caching, headroom, the strength models and the transfer study ran on an RTX 4050 laptop (PyTorch 2.5.1, CUDA 12.1), both with 6 GB of memory. The code also runs on Apple silicon and the CPU; the CPU path was verified. Timings are compared only on one machine.

The application \texttt{experiment\_app.py} has seven tabs in three tiers: the interface, the intervention engine and the model and data layer (cached GPT-2, SAEs, fact sets, per-prompt caches). \emph{Hybrid mute and boost} carries both editors behind a choice of editing method (the fixed sweep with its strengths, cumulative sweep, pool refill, safety filter and GPU-batched switch, or gradient descent with its steps, learning rate, margin and weights), with a live tracker, a stopwatch and a session-history record. \emph{Prototype lab} runs gradient descent with the optional additive edit, shows baseline and edited text side by side and applies the same edit to follow-up prompts. \emph{Batch} runs studies in background workers; \emph{Repair \& SAE limit} shows the diagnostics of Section~\ref{sec:repair}; \emph{Sequential vs Batched Proof} times the two code paths; \emph{Monosemanticity Analysis} inspects features; \emph{Session history} stores every run. The edit lives in \texttt{src/hooks.py} (Equation~\ref{eq:delta}), \texttt{src/editing.py} (screening, causal ranking and the safety checks of Equation~\ref{eq:safe}), \texttt{src/hybrid\_runner.py} (the headless sweep, which reproduces the application exactly) and \texttt{src/gradient\_editing.py} (Equations~\ref{eq:a} to~\ref{eq:add}). The transfer study adds only new files (\texttt{src/transfer/}, tools \texttt{00} to \texttt{27} in \texttt{tools/transfer/}); older functions are imported and never edited, so earlier results cannot change.

\section{Correctness Checks and Engineering Lessons}

Every fast path was checked against a slower reference before use, and the checks are kept as scripts: batched against sequential ranking (identical feature lists, probabilities to $10^{-9}$, 60 of 60 safety decisions); the cached-state path against the full model (logits identical on one machine, within $10^{-5}$ on the other); a random per-feature edit through the cache path against the real hook (same target rank); and the gradient-descent tuning path against the real hook in every recorded run. Built-in sanity checks stop a run on a mismatch; one caught a cache-builder error (a multiplier 1.0 passed where 0.0 was meant) before any data was written. Three problems could have invalidated results without any error message.
\begin{itemize}
  \item \textbf{A thread race removed the start-of-text token.} Tokenising with the start token switched off temporarily changes a setting shared by the whole model in TransformerLens; when two Streamlit sessions overlapped, the wrong value could be written back permanently and every later prompt was tokenised without its start token (rank 28 instead of 22 on one prompt). Of 57 recorded application runs, 16 had wrong baselines, all from one server process, and none was used in a table. The fix tokenises directly and never touches the model's configuration.
  \item \textbf{A timer measured the network} (46 s against about 10 s of model time), and \textbf{chained hooks disagreed with the combined check} (rank 66 against 59); look-ups moved out of the timer and one combined hook is used.
  \item \textbf{GPU memory.} A padded training set of 7,550 texts (3.8 GB) nearly filled the 6 GB card and made runs five to ten times slower without any error (6.2 hours instead of one). Training tensors now stay on the CPU.
\end{itemize}

\section{Reproducibility and Development Process}

Every tool prints its pass rule, saves its raw output as a log and writes a result file with per-run details. Logs and results of every experiment are copied into a data pack under version control, and the figures of this report are drawn from the pack by scripts that run no model; larger local artefacts (2.6 GB) are archived with a manifest of sizes and checksums. Optimisation followed a \emph{one-bounded-step} rule: identify one bottleneck, measure the reference, make one change, verify equivalence on a fixed case, log, and only then approve the next step.